\section{三类学习范式对比}

在正式进入强化学习之前，有必要将其与前面学过的两种学习范式进行系统比较，以便理解其独特定位。

\begin{figure}[H]
\centering
\includegraphics[width=0.95\textwidth]{slides-images/slide-07.jpg}
\caption{三类学习问题的对比：Supervised Learning、Unsupervised Learning 与 Reinforcement Learning\protect\footnotemark}
\end{figure}
\footnotetext{Slides 第7页。}

\subsection{Supervised Learning（监督学习）}

\begin{itemize}
    \item \textbf{数据}：成对的 $(x, y)$，其中 $x$ 是输入数据，$y$ 是标签
    \item \textbf{目标}：学习一个从 $x$ 到 $y$ 的映射函数
    \item \textbf{类比}：给模型展示大量苹果的图片，并告诉它"这是苹果"，模型学会识别新的苹果图片
\end{itemize}

\subsection{Unsupervised Learning（无监督学习）}

\begin{itemize}
    \item \textbf{数据}：只有 $x$，没有标签
    \item \textbf{目标}：学习数据的内在结构和分布
    \item \textbf{类比}：给模型展示大量相似物体的图片（不告诉它是什么），模型学会发现它们的共同特征
\end{itemize}

\subsection{Reinforcement Learning（强化学习）}

\begin{itemize}
    \item \textbf{数据}：State-Action pairs（状态-动作对）
    \item \textbf{目标}：最大化在多个时间步上的\textbf{未来累积奖励}
    \item \textbf{类比}：不需要告诉模型"这是苹果"，而是让它通过交互学会"吃苹果可以获得营养从而存活"
\end{itemize}

\begin{importantbox}{三种范式的本质区别}
\begin{itemize}
    \item \textbf{Supervised Learning}：从\textbf{标注数据}中学习输入到输出的映射
    \item \textbf{Unsupervised Learning}：从\textbf{无标注数据}中发现隐藏结构
    \item \textbf{Reinforcement Learning}：通过与\textbf{环境的交互}学习最优决策策略，目标是最大化长期累积回报
\end{itemize}
学习算法（Algorithm）是与网络架构（Architecture）无关的（agnostic）。同一种架构（如 CNN）可以用于不同的学习范式。
\end{importantbox}

\subsection{本章小结}
三种学习范式分别对应不同的数据形式和学习目标。强化学习的独特之处在于：（1）数据不是预先收集的，而是通过交互动态生成的；（2）目标不是预测标签或发现结构，而是学习最优行为策略以最大化长期回报；（3）当前的决策会影响后续能观察到的数据（非 i.i.d.）。


